\documentclass[10pt,a4paper]{amsart}
\usepackage{amsmath,amssymb,amsthm,graphicx}
\usepackage{stmaryrd}
\newtheorem{theorem}{Theorem}
\newtheorem{definition}[theorem]{Definition}
\newtheorem{problem}[theorem]{Problem}
\usepackage{hyperref}
\providecommand{\vmathbb}[1]{\mathbb{#1}}
\NewCommandCopy{\oldGamma}{\Gamma}\renewcommand{\Gamma}{{\mathit{\oldGamma}}}
\NewCommandCopy{\oldDelta}{\Delta}\renewcommand{\Delta}{{\mathit{\oldDelta}}}
\NewCommandCopy{\oldTheta}{\Theta}\renewcommand{\Theta}{{\mathit{\oldTheta}}}
\newcommand{\Ty}{\mathsf{Ty}}
\newcommand{\Tm}{\mathsf{Tm}}
\newcommand{\Tys}{\mathsf{Tys}}
\newcommand{\Tms}{\mathsf{Tms}}
\newcommand{\Set}{\mathsf{Set}}
\newcommand{\Prop}{\mathsf{Prop}}
\newcommand{\U}{\mathsf{U}}
\newcommand{\El}{\mathsf{El}}
\newcommand{\Up}{\mathsf{U}^{+}}
\newcommand{\ra}{\rightarrow}
\newcommand{\Ra}{\Rightarrow}
\newcommand{\blank}{\mathord{\hspace{1pt}\text{--}\hspace{1pt}}} %from the book
\renewcommand{\iff}{\leftrightarrow}
\newcommand{\lam}{\mathsf{lam}}
\newcommand{\app}{\mathsf{app}}
\newcommand{\el}{\mathsf{el}}
\newcommand{\Con}{\mathsf{Con}}
\newcommand{\Sub}{\mathsf{Sub}}
\newcommand{\ext}{\operatorname{\triangleright}}
\newcommand{\blackext}{\operatorname{\blacktriangleright}}
\newcommand{\p}{\mathsf{p}}
\newcommand{\q}{\mathsf{q}}
\newcommand{\id}{\mathsf{id}}
\newcommand{\For}{\mathsf{For}}
\newcommand{\Pf}{\mathsf{Pf}}
\newcommand{\Eq}{\mathsf{Eq}}
\newcommand{\Id}{\mathsf{Id}}
\newcommand{\Lam}{\mathsf{Lam}}
\newcommand{\LAM}{\mathsf{LAM}}
\newcommand{\Kind}{\mathsf{Kind}}
\newcommand{\Val}{\mathsf{Val}}
\newcommand{\Srt}{\mathsf{Srt}}
\newcommand{\bind}{\mathbin{\gg\!\!=}}
\newcommand{\fst}{\mathsf{fst}}
\newcommand{\snd}{\mathsf{snd}}
\newcommand{\C}{\mathcal{C}}
\newcommand{\D}{\mathcal{D}}
\newcommand{\E}{\mathcal{E}}
\newcommand{\PSh}{\mathsf{PSh}}
\renewcommand{\tt}{\mathsf{tt}}
\newcommand{\y}{\mathsf{y}}
\newcommand{\yl}{\mathsf{yl}}
\newcommand{\code}{\mathsf{c}}
\newcommand{\cur}{\urcorner}
\newcommand{\cul}{\ulcorner}
\title{Naive and direct first-order semantics of second-order generalised algebraic signatures yield isomorphic models}
\author{Ambrus Kaposi, Szumi Xie}
\date{doi:10.4230/LIPIcs.FSCD.2024.10; CC BY4.0 (official publisher paper; selected published text only)}
\begin{document}
\maketitle
\noindent\textit{Editorial note.} Source https://drops.dagstuhl.de/entities/document/10.4230/LIPIcs.FSCD.2024.10. Copyright Ambrus Kaposi and Szumi Xie, CC BY4.0. This item contains published Theorem22 and complete author proof; native author TeX is used only for text matching the licensed published article. Theorem27 and its supplementary detailed construction are outside this item.\par\medskip
\noindent\textit{Allowed standard prerequisites (editorial).}\begin{itemize}
\item Categories with families, presheaves, Yoneda lemma and CwF pseudomorphisms. Publisher references19,31,37. These are standard prerequisites cited by the author. The actual ToS/ToS+ and interpretation constructions are retained in the editable material.
\end{itemize}\medskip
\setcounter{theorem}{11}
universe (a type and a family over it), $\top$ and $\Sigma$ types,
equality with reflection, and a $\Pi$ type with $\U$-domain.
\pagebreak
\begin{definition}[ToS: the theory of GAT signatures]\label{def:gat}
  \begin{alignat*}{10}
    & \Ty && : \Set                                             && \Sigma && :(A:\Ty)\ra(\Tm\,A\ra\Ty)\ra\Ty \\
    & \Tm && : \Ty\ra\Set                                       && (\blank,\blank) && : (a:\Tm\,A)\times\Tm\,(B\,a)\cong\Tm\,(\Sigma\,A\,B) : \mathsf{fst},\mathsf{snd} \\         
    & \U  && : \Ty                                              && \Pi && : (a:\Tm\,\U)\ra(\Tm\,(\El\,a)\ra\Ty)\ra\Ty \\                                                 
    & \El && : \Tm\,\U\ra\Ty                                    && \lam && : \big((x:\Tm\,(\El\,a))\ra\Tm\,(B\,x)\big)\cong\Tm\,(\Pi\,a\,B) : \blank\cdot\blank \\     
    & \top && : \Ty                                             && \Eq && : (A:\Ty)\ra\,\Tm\,A\ra\,\Tm\,A\ra\Ty \\                                                                 
    & \tt && : \vmathbb{1} \cong \Tm\,\top \hspace{6em} && \mathsf{refl} && : (u = v)\cong\Tm\,(\Eq\,A\,u\,v) : \mathsf{reflect}                                    
  \end{alignat*}
\end{definition}

\setcounter{theorem}{12}
For the theory of SOGAT signatures (ToS$^+$), we add a new universe $\U^+$ of sorts for which
variables are allowed: with the help of these we can write second
order functions. $\U^+$ is a subuniverse of $\U$ (witnessed by
$\el^+$) and has a $\Pi$ type with $\U^+$-domain and $\U$-codomain.
\begin{definition}[ToS$^+$: the theory of SOGAT signatures]\label{def:sogat}
  We extend ToS with the following.
  \begin{alignat*}{10}
    & \U^+ && : \Ty                                   && \pi^+ && : (a^+ : \Tm\,\U^+)\ra\Big(\Tm\,\big(\El\,(\el^+\,a^+)\big)\ra\Tm\,\U\Big)\ra\Tm\,\U \\                                          
    & \el^+ && : \Tm\,\U^+ \ra \Tm\,\U \hspace{1.7em} && \lam^+ && : \big(x:\El\,(\el^+\,a^+)\big)\ra\Tm\,(\El\,(b\,x))\cong\Tm\,\big(\El\,(\pi^+\,a^+\,b)\big) : \blank\cdot^+\blank    
  \end{alignat*}
\end{definition}
The first-order version is \autoref{def:sogat_gat}. A (presentation of
a) GAT is defined as a closed type in the syntax of ToS$^+$.

\setcounter{theorem}{13}
\section{Naive semantics of SOGAT signatures}
\label{sec:naive}

In this section, for any SOGAT signature, we define a notion of
first-order model. The idea is that a model is a category together
with the presheaf interpretation of the signature over that category:
the category of presheaves supports a universe, $\Pi$ types, and so
on, so we directly use these when interpreting the type formers of ToS$^+$.
We assume basic working knowledge of
categories with families (CwFs
[19]).
\begin{definition}[First-order model of ToS]\label{def:gat_gat}
  A first-order model of ToS is a CwF (sorts
  are denoted $\Con$, $\Sub$, $\Ty$, $\Tm$, the empty context is $\diamond$, the empty substitution is $\epsilon:\Sub\,\Gamma\,\diamond$, context extension is
  $\blank\ext\blank : (\Gamma:\Con)\ra\Ty\,\Gamma\ra\Con$ with $\p\circ\blank,\q[\blank] :
  \Sub\,\Delta\,(\Gamma\ext A) \cong
  (\gamma:\Sub\,\Delta\,\Gamma)\times\Tm\,\Delta\,(A[\gamma]) :
  \blank\mathbin{{,}}\blank$) equipped with:
  \begin{itemize}
  \item $\top$ and $\Sigma$ types given by isomorphisms \\
    $
    \tt : \vmathbb{1} \cong \Tm\,\Gamma\,\top, \hspace{4em} (\blank,\blank) : (a:\Tm\,\Gamma\,A)\times\Tm\,\Gamma\,(B[\id,a]) \cong \Tm\,\Gamma\,(\Sigma\,A\,B) : \fst,\snd.
    $
  \item A universe given by $\U : \Ty\,\Gamma$ and $\El : \Tm\,\Gamma\,\U\ra\Ty\,\Gamma$.
  \item A function space with domain in $\U$, that is
    $\Pi:(a:\Tm\,\Gamma\,\U)\ra\Ty\,(\Gamma\ext\El\,a)\ra\Ty\,\Gamma$,
    with an isomorphism $\lam : \Tm\,(\Gamma\ext\El\,a)\,B \cong 
    \Tm\,\Gamma\,(\Pi\,a\,B) : \app$.
  \item A strict equality type $\Eq$ with reflection and uniqueness of
    identity proofs.
  \item All the operations listed above are natural in $\Gamma$.
  \end{itemize}
\end{definition}
\begin{definition}[First-order model of ToS$^+$]\label{def:sogat_gat}
  A first-order model of ToS$^+$ is a
  first-order model of ToS equipped with:
  \begin{itemize}
  \item Another universe $\U^+ : \Ty\,\Gamma$ that is a subuniverse of
    $\U$ i.e.\ $\el^+:\Tm\,\Gamma\,\U^+\ra\Tm\,\Gamma\,\U$.
  \item $\U$ is closed under functions with $\U^+$-domain, i.e.\ 
    $\pi^+ :
    (a^+:\Tm\,\Gamma\,\U^+)\ra\Tm\,\big(\Gamma\ext\El\,(\el^+\,a^+)\big)\,\U\ra\Tm\,\Gamma\,\U$
    with $\lam^+ :
    \Tm\,\big(\Gamma\ext\El\,(\el^+\,a)\big)\,(\El\,b) \cong 
    \Tm\,\Gamma\,\big(\El\,(\pi^+\,a^+\,b)\big) : \app^+$.
  \item All the operations listed above are natural in $\Gamma$.
  \end{itemize}
\end{definition}
\begin{problem}[$\PSh$]\label{prob:presheaf}
  Presheaves over a category $C$ form a CwF equipped with $\top$, $\Sigma$ types,
  an equality type with reflection, $\Pi$ types and a Coquand-universe
  $\U$ with $\code:\Ty\,\Gamma\cong\Tm\,\Gamma\,\U:\El$. Unlike in
  original Definition 11, we omit writing universe indices for
  readibility.
\end{problem}
\begin{proof}[Construction]
  We recall the main parts of the construction
  [31] for fixing notations.  $\Gamma:\Con$ is a
  presheaf, that is a family of sets $\Gamma:\C\ra\Set$ with
  reindexing $\gamma_I[f]_\Gamma : \Gamma\,J$ for $\gamma_I:\Gamma\,I$
  and $f:\C(J,I)$ such that $\gamma_I[f\circ g]_\Gamma =
  \gamma_I[f]_\Gamma[g]_\Gamma$ and $\gamma_I[\id]_\Gamma =
  \gamma_I$. A $\sigma:\Sub\,\Delta\,\Gamma$ is a function $\sigma :
  \Delta\,I\ra\Gamma\,I$ such that $(\sigma\,\delta_I)[f]_\Gamma =
  \sigma\,(\delta_I[f]_\Delta)$. A type $A:\Ty\,\Gamma$ is a dependent
  presheaf containing a family $A : (I:\C)\ra\Gamma\,I\ra\Set$ with
  reindexing $a_I[f]_A : A\,J\,(\gamma_I[f]_\Gamma)$ for
  $a_I:A\,I\,\gamma_I$ and $f:\C(J,I)$ satisfying functoriality. Type
  substitution is $A[\gamma]\,I\,\delta_I :=
  A\,I\,(\gamma\,\delta_I)$. A term $a:\Tm\,\Gamma\,A$ is a function
  $a : (\gamma_I:\Gamma\,I)\ra A\,I\,\gamma_I$ such that
  $(a\,\gamma_I)[f]_A = a\,(\gamma_I[f]_\Gamma)$. Term substitution is
  $a[\gamma]\,\delta_I := a\,(\gamma\,\delta_I)$. The empty context is
  constant unit: $\diamond\,I := \vmathbb{1}$. Context extension is
  pointwise: $(\Gamma\ext A)\,I := (\gamma_I:\Gamma\,I)\times
  A\,I\,\gamma_I$, its universal property is given by projections and
  pairing for metatheoretic $\Sigma$ types. $\top$, $\Sigma$ and $\Eq$
  are pointwise. We have the functor Yoneda $\y : \C\ra\PSh(\C)$
  defined by $\y\,I\,J := \C(J,I)$, and we use this to define the
  universe by $\U\,I\,\gamma_I := \Ty\,(\y\,I)$. We observe that
  $\gamma_I[\blank]_\Gamma : \Sub\,(\y\,I)\,\Gamma$ (forward part of
  Yoneda lemma), and define $\Pi\,A\,B\,I\,\gamma_I := \Tm\,(\y\,I\ext
  A[\gamma_I[\blank]_\Gamma])\,(B[\gamma_I[\blank]_\Gamma\circ\p,\q])$.
\end{proof}
\begin{problem}[Locally representable types]
  The CwF of presheaves can be extended to a CwF$^+$, which means a
  CwF with a subsort of $\Ty$ called $\Ty^+$ and a $\Pi^+$ type with
  domain in $\Ty^+$, i.e.\ $\Pi^+ :
  (A:\Ty^+\,\Gamma)\ra\Ty\,(\Gamma\ext A)\ra\Ty\,\Gamma$ with $\lam^+
  : \Tm\,(\Gamma\ext A)\,B\cong\Tm\,\Gamma\,(\Pi^+\,A\,B) : \app^+$,
  natural in $\Gamma$. $\Ty^+$ is classified by the Coquand universe
  $\U^+$.
\end{problem}
\begin{proof}[Construction]
  An element $A : \Ty^+\,\Gamma$ is an $A : \Ty\,\Gamma$ together with
  $\blank\ext_A\blank : (I:\C)\ra\Gamma\,I\ra\C$ and an isomorphism
  $\p_A\circ\blank,\q_A[\blank]_A : \C(J,I\ext_A\gamma_I) \cong
  (f:\C(J,I))\times A\,J\,(\gamma_I[f]_\Gamma) :
  \blank\mathbin{,_A}\blank$ natural in $J$. So
  $\p_A:\C(I\ext_A\gamma_I,I)$ and
  $\q_A:A\,(I\ext_A\gamma_I)\,(\gamma_I[\p_A]_\Gamma)$. Substitution
  is given by $I\ext_{A[\gamma]}\delta_I := I\ext_A\gamma\,\delta_I$
  and we have $\C(J,I\ext_{A[\gamma]}\delta_I) =
  \C(J,I\ext_A\gamma\,\delta_I) \cong (f:\C(J,I))\times
  A\,J\,(\gamma\,\delta_I[f]_\Gamma) = (f:\C(J,I))\times
  A[\gamma]\,J\,(\delta_I[f]_\Delta)$. We define $\Pi^+$ using the
  $\ext_A$ operator which comes with $A$,
  i.e.\ $\Pi^+\,A\,B\,I\,\gamma_I :=
  B\,(I\ext_A\gamma_I)\,(\gamma_I[\p_A]_\Gamma, \q_A)$,
  $b_{I'}[f]_{\Pi^+\,A\,B} := b_{I'}[f\circ\p_A\mathbin{,_A}\q_A]$,
  $\lam^+\,b\,\gamma_I := b\,(\gamma_I[\p_A]_\Gamma,\q_A)$ and
  $\app^+\,t\,(\gamma_I,a_I) :=
  (t\,\gamma_I)[\id_I\mathbin{,_A}a_I]_B$. Like $\U$,
  $\U^+\,I\,\gamma_I := \Ty^+\,(\y\,I)$.
\end{proof}
\begin{definition}[Naive semantics]\label{def:naive}
  Given a category $\C$, $\PSh(\C)$ is a model of ToS$^+$ choosing $\U := \U$, $\El\,a := \El\,a$, $\Pi\,a\,B :=
  \Pi\,(\El\,a)\,B$, $\U^+ := \U^+$, $\el^+\,a^+ :=
  \code\,(\El^+\,a^+)$, $\pi^+\,a^+\,b :=
  \code\,\big(\Pi^+\,(\El^+\,a^+)\,(\El\,b)\big)$. Recall that a SOGAT
  signature $\Omega$ is an element of $\Ty\,\diamond$ in the syntax of ToS$^+$.
  A naive model of $\Omega$ is a category with a
  terminal object $\diamond$ together with the interpretation of $\Omega$ in
  presheaves over this category, i.e.\ 
  $(\C:\mathsf{Cat}^\diamond)\times\Tm_{\PSh(\C)}\,\diamond\llbracket\Omega
  \rrbracket_{\PSh(\C)}$.
\end{definition}

\setcounter{theorem}{18}
\section{Direct semantics of SOGAT signatures}
\label{sec:warmup}

In this section, we define first-order models of SOGATs using a more
careful version of the presheaf model. We make sure that no
Yoneda-encodings are present in the semantics using the idea of
two-level type theory
[4, 10]
where presheaves over a CwF include a universe of ``inner types''
coming from the CwF. We extend two-level type theory with a separate function space
where the domain is an inner type. This function space is isomorphic
to the usual presheaf function space, but has a simpler semantics.
\begin{problem}[Presheaves over a CwF]\label{prob:psh_gat}
  If $\C$ is a CwF, then $\PSh(\C)$ models ToS without using the usual
  presheaf $\U$ and $\Pi$.
\end{problem}
\begin{proof}[Construction]
  We interpret $\top$, $\Sigma$, $\Eq$ as in \autoref{prob:presheaf},
  but define $\U$, $\El$ and $\Pi$ by $\Ty_\C$, $\Tm_\C$ and
  $\ext_\C$, respectively: $\U\,I\,\gamma_I := \Ty_\C\,I$,
  
  $\El\,a\,I\,\gamma_I := \Tm_\C\,I\,(a\,\gamma_I)$,
  
  $\Pi\,a\,B\,I\,\gamma_I := B\,(I\ext_\C a\,\gamma_I)\,(\gamma_I[\p_\C]_\Gamma,\q_\C)$ with
  $\lam\,b\,\gamma_I := b\,(\gamma_I[\p_\C]_\Gamma,\q_C)$ and
  $\app\,t\,(\gamma_I,a_I) := t\,\gamma_I[\id_I\mathbin{,_\C}a_I]_B$.
\end{proof}
\begin{problem}[Presheaves over a CwF$^+$]
  If the category $\C$ is a CwF$^+$, then the previous model extends
  to a model of ToS$^+$ (\autoref{def:sogat_gat}).
\end{problem}
\begin{proof}[Construction]
  We interpret $\U^+$, $\el^+$ and $\pi^+$ by $\Ty^+_\C$, identity and
  $\Pi^+_\C$, respectively: $\U^+\,I\,\gamma_I := \Ty^+_\C\,I$,
  $\el^+\,a\,\gamma_I := a\,\gamma_I$, $\pi^+\,a\,b\,\gamma_I :=
  \Pi^+\,(a\,\gamma_I)\,\big(b\,(\gamma_I[\p_\C]_\Gamma,\q_\C)\big)$,
  $\lam^+\,t\,\gamma_I :=
  \lam^+_\C\,\big(t\,(\gamma_I[\p_\C]_\Gamma,\q_\C)\big)$,
  $\app^+\,t\,(\gamma_I,a_I) :=
  \app^+_\C\,(t\,\gamma_I)[\id_I\mathbin{,_\C}a_I]_{\Tm_\C}$.
\end{proof}
\begin{definition}[Direct semantics]\label{def:direct}
  A direct model of a SOGAT signature $\Omega$ is a category $\C$ with a
  terminal object together with the interpretation of $\Omega$ in
  presheaves over presheaves over $\C$, evaluated at the terminal
  presheaf: $(\C:\mathsf{Cat}^\diamond) \times
  \llbracket\Omega\rrbracket_{\PSh(\PSh(\C))}\,\diamond_{\PSh(\C)}\,{\star}$. Note
  that this makes sense because
  $\PSh(\C):\mathsf{CwF}^+$, hence $\PSh\big(\PSh(\C)\big)$ is a model
  of ToS$^+$.
\end{definition}

\setcounter{theorem}{21}
\begin{theorem}\label{thm:naive_direct}
  For any signature, the naive and direct semantics result in
  isomorphic notions of models.
\end{theorem}

\setcounter{theorem}{21}
\begin{proof}
  We fix a $\C:\mathsf{Cat}^\diamond$, and denote $\D := \PSh(\C)$ and
  $\E := \PSh(\D)$. $\D$ is a model of ToS$^+$
  via \autoref{def:naive} and $\E$ is a model via
  \autoref{def:direct}, and Yoneda navigates between them (it is not
  only a functor, but a CwF pseudomorphism
  [37]). By induction on the syntax of ToS$^+$, we define $\alpha$ for contexts,
  substitutions, types and terms: 
  $\alpha_\Gamma : \Sub_{\E}\,\llbracket\Gamma\rrbracket_\E\,(\y\,\llbracket\Gamma\rrbracket_\D)$,
  $\alpha_\gamma : \alpha_\Gamma\circ\llbracket\gamma\rrbracket_\E = \y\,\llbracket\gamma\rrbracket_\D\circ\alpha_\Delta$,
  $\alpha_A : \llbracket A\rrbracket_\E \cong \y\,\llbracket A\rrbracket_\D[\alpha_\Gamma]$,
  $\alpha_a : \alpha_A[\id,\llbracket a\rrbracket_\E] = \y\,\llbracket a\rrbracket_\D[\alpha_\Gamma]$.
  




  For a signature $\Omega:\Ty\,\diamond$, we thus obtain $\llbracket\Omega\rrbracket_\E\,\diamond_\D\,{\star} \cong \y\,\llbracket\Omega\rrbracket_\D[\alpha_\Gamma]\,\diamond_\D\,{\star} = \Tm_\D\,\diamond_\D\,\llbracket\Omega\rrbracket_\D$.
\end{proof}
Note that there is no size issue when stating the isomorphism because
even if $\E$ is one level up compared to $\D$, we only use small
components from $\E$ when evaluating into it.

\begin{thebibliography}{99}
\bibitem[1]{source1} Benedikt Ahrens, Jacopo Emmenegger, Paige Randall North, and Egbert Rijke. B-systems and C-systems are equivalent. The Journal of Symbolic Logic, pages 1-9, 2023. URL: https://doi.org/10.1017/jsl.2023.41 .
\bibitem[2]{source2} Benedikt Ahrens, André Hirschowitz, Ambroise Lafont, and Marco Maggesi. Modular specification of monads through higher-order presentations. In Herman Geuvers, editor, 4th International Conference on Formal Structures for Computation and Deduction, FSCD 2019, June 24-30, 2019, Dortmund, Germany, volume 131 of LIPIcs, pages 6:1-6:19. Schloss Dagstuhl - Leibniz-Zentrum für Informatik, 2019. URL: https://doi.org/10.4230/LIPICS.FSCD.2019.6 .
\bibitem[3]{source3} Guillaume Allais, Robert Atkey, James Chapman, Conor McBride, and James McKinna. A type- and scope-safe universe of syntaxes with binding: their semantics and proofs. J. Funct. Program., 31:e22, 2021. URL: https://doi.org/10.1017/S0956796820000076 .
\bibitem[4]{source4} Thorsten Altenkirch, Paolo Capriotti, and Nicolai Kraus. Extending homotopy type theory with strict equality. In Jean-Marc Talbot and Laurent Regnier, editors, 25th EACSL Annual Conference on Computer Science Logic, CSL 2016, August 29 - September 1, 2016, Marseille, France, volume 62 of LIPIcs, pages 21:1-21:17. Schloss Dagstuhl - Leibniz-Zentrum für Informatik, 2016. URL: https://doi.org/10.4230/LIPICS.CSL.2016.21 .
\bibitem[5]{source5} Thorsten Altenkirch, Yorgo Chamoun, Ambrus Kaposi, and Michael Shulman. Internal parametricity, without an interval. Proc. ACM Program. Lang., 8(POPL):2340-2369, 2024. URL: https://doi.org/10.1145/3632920 .
\bibitem[6]{source6} Thorsten Altenkirch and Ambrus Kaposi. Normalisation by evaluation for dependent types. In Delia Kesner and Brigitte Pientka, editors, 1st International Conference on Formal Structures for Computation and Deduction, FSCD 2016, June 22-26, 2016, Porto, Portugal, volume 52 of LIPIcs, pages 6:1-6:16. Schloss Dagstuhl - Leibniz-Zentrum für Informatik, 2016. URL: https://doi.org/10.4230/LIPICS.FSCD.2016.6 .
\bibitem[7]{source7} Thorsten Altenkirch and Ambrus Kaposi. Type theory in type theory using quotient inductive types. In Rastislav Bodík and Rupak Majumdar, editors, Proceedings of the 43rd Annual ACM SIGPLAN-SIGACT Symposium on Principles of Programming Languages, POPL 2016, St. Petersburg, FL, USA, January 20 - 22, 2016, pages 18-29. ACM, 2016. URL: https://doi.org/10.1145/2837614.2837638 .
\bibitem[8]{source8} Thorsten Altenkirch, Ambrus Kaposi, Artjoms Sinkarovs, and Tamás Végh. Combinatory logic and lambda calculus are equal, algebraically. In Marco Gaboardi and Femke van Raamsdonk, editors, 8th International Conference on Formal Structures for Computation and Deduction, FSCD 2023, July 3-6, 2023, Rome, Italy, volume 260 of LIPIcs, pages 24:1-24:19. Schloss Dagstuhl - Leibniz-Zentrum für Informatik, 2023. URL: https://doi.org/10.4230/LIPICS.FSCD.2023.24 .
\bibitem[9]{source9} Thorsten Altenkirch and Bernhard Reus. Monadic presentations of lambda terms using generalized inductive types. In Jörg Flum and Mario Rodríguez-Artalejo, editors, Computer Science Logic, 13th International Workshop, CSL '99, 8th Annual Conference of the EACSL, Madrid, Spain, September 20-25, 1999, Proceedings, volume 1683 of Lecture Notes in Computer Science, pages 453-468. Springer, 1999. URL: https://doi.org/10.1007/3-540-48168-0\_32 .
\bibitem[10]{source10} Danil Annenkov, Paolo Capriotti, Nicolai Kraus, and Christian Sattler. Two-level type theory and applications. Mathematical Structures in Computer Science, 33(8):688-743, 2023. URL: https://doi.org/10.1017/S0960129523000130 .
\bibitem[11]{source11} Samy Avrillon. Logic as a second-order generalized algebraic theory, 2023. Report on the 3-month research internship at the Faculty of Informatics of ELTE. URL: https://github.com/MysaaJava/m1-internship/releases/download/project-report/Avrillon-02.pdf .
\bibitem[12]{source12} Henk Barendregt. Introduction to generalized type systems. Journal of Functional Programming, 1(2):125-154, 1991. URL: https://doi.org/10.1017/S0956796800020025 .
\bibitem[13]{source13} Andrej Bauer, Philipp G. Haselwarter, and Peter LeFanu Lumsdaine. A general definition of dependent type theories. CoRR, abs/2009.05539, 2020. URL: https://arxiv.org/abs/2009.05539 .
\bibitem[14]{source14} Rafaël Bocquet. External univalence for second-order generalized algebraic theories. CoRR, abs/2211.07487, 2022. URL: https://doi.org/10.48550/arXiv.2211.07487 .
\bibitem[15]{source15} Rafaël Bocquet. Towards coherence theorems for equational extensions of type theories. CoRR, abs/2304.10343, 2023. URL: https://doi.org/10.48550/arXiv.2304.10343 .
\bibitem[16]{source16} Rafaël Bocquet, Ambrus Kaposi, and Christian Sattler. For the metatheory of type theory, internal sconing is enough. In Marco Gaboardi and Femke van Raamsdonk, editors, 8th International Conference on Formal Structures for Computation and Deduction, FSCD 2023, July 3-6, 2023, Rome, Italy, volume 260 of LIPIcs, pages 18:1-18:23. Schloss Dagstuhl - Leibniz-Zentrum für Informatik, 2023. URL: https://doi.org/10.4230/LIPICS.FSCD.2023.18 .
\bibitem[17]{source17} Paolo Capriotti. Notions of type formers. In Ambrus Kaposi, editor, 23rd International Conference on Types for Proofs and Programs, TYPES 2017. Eötvös Loránd University, 2017. URL: http://types2017.elte.hu/proc.pdf\#page=77 .
\bibitem[18]{source18} John Cartmell. Generalised algebraic theories and contextual categories. Ann. Pure Appl. Log., 32:209-243, 1986. URL: https://doi.org/10.1016/0168-0072(86)90053-9 .
\bibitem[19]{source19} Simon Castellan, Pierre Clairambault, and Peter Dybjer. Categories with families: Unityped, simply typed, and dependently typed. CoRR, abs/1904.00827, 2019. URL: https://arxiv.org/abs/1904.00827 .
\bibitem[20]{source20} Thierry Coquand. Canonicity and normalization for dependent type theory. Theor. Comput. Sci., 777:184-191, 2019. URL: https://doi.org/10.1016/J.TCS.2019.01.015 .
\bibitem[21]{source21} Marcelo P. Fiore and Chung-Kil Hur. Second-order equational logic (extended abstract). In Anuj Dawar and Helmut Veith, editors, Computer Science Logic, 24th International Workshop, CSL 2010, 19th Annual Conference of the EACSL, Brno, Czech Republic, August 23-27, 2010. Proceedings, volume 6247 of Lecture Notes in Computer Science, pages 320-335. Springer, 2010. URL: https://doi.org/10.1007/978-3-642-15205-4\_26 .
\bibitem[22]{source22} Marcelo P. Fiore, Andrew M. Pitts, and S. C. Steenkamp. Quotients, inductive types, and quotient inductive types. Log. Methods Comput. Sci., 18(2), 2022. URL: https://doi.org/10.46298/LMCS-18(2:15)2022 .
\bibitem[23]{source23} Marcelo P. Fiore, Gordon D. Plotkin, and Daniele Turi. Abstract syntax and variable binding. In 14th Annual IEEE Symposium on Logic in Computer Science, Trento, Italy, July 2-5, 1999, pages 193-202. IEEE Computer Society, 1999. URL: https://doi.org/10.1109/LICS.1999.782615 .
\bibitem[24]{source24} Jonas Frey. Duality for clans: a refinement of gabriel-ulmer duality. CoRR, abs/2308.11967, 2023. URL: https://doi.org/10.48550/arXiv.2308.11967 .
\bibitem[25]{source25} Gaëtan Gilbert, Jesper Cockx, Matthieu Sozeau, and Nicolas Tabareau. Definitional proof-irrelevance without K. Proc. ACM Program. Lang., 3(POPL):3:1-3:28, 2019. URL: https://doi.org/10.1145/3290316 .
\bibitem[26]{source26} Daniel Gratzer. Normalization for multimodal type theory. In Christel Baier and Dana Fisman, editors, LICS '22: 37th Annual ACM/IEEE Symposium on Logic in Computer Science, Haifa, Israel, August 2 - 5, 2022, pages 2:1-2:13. ACM, 2022. URL: https://doi.org/10.1145/3531130.3532398 .
\bibitem[27]{source27} Robert Harper. Practical Foundations for Programming Languages (2nd. Ed.). Cambridge University Press, 2016. URL: https://www.cs.cmu.edu/\%7Erwh/pfpl/index.html .
\bibitem[28]{source28} Robert Harper. An equational logical framework for type theories. CoRR, abs/2106.01484, 2021. URL: https://arxiv.org/abs/2106.01484 .
\bibitem[29]{source29} Robert Harper, Furio Honsell, and Gordon D. Plotkin. A framework for defining logics. J. ACM, 40(1):143-184, 1993. URL: https://doi.org/10.1145/138027.138060 .
\bibitem[30]{source30} Philipp G. Haselwarter and Andrej Bauer. Finitary type theories with and without contexts. J. Autom. Reason., 67(4):36, 2023. URL: https://doi.org/10.1007/S10817-023-09678-Y .
\bibitem[31]{source31} Martin Hofmann. Syntax and semantics of dependent types. In Semantics and Logics of Computation, pages 79-130. Cambridge University Press, 1997.
\bibitem[32]{source32} Martin Hofmann. Semantical analysis of higher-order abstract syntax. In 14th Annual IEEE Symposium on Logic in Computer Science, Trento, Italy, July 2-5, 1999, pages 204-213. IEEE Computer Society, 1999. URL: https://doi.org/10.1109/LICS.1999.782616 .
\bibitem[33]{source33} Jasper Hugunin. Why not W? In Ugo de'Liguoro, Stefano Berardi, and Thorsten Altenkirch, editors, 26th International Conference on Types for Proofs and Programs, TYPES 2020, March 2-5, 2020, University of Turin, Italy, volume 188 of LIPIcs, pages 8:1-8:9. Schloss Dagstuhl - Leibniz-Zentrum für Informatik, 2020. URL: https://doi.org/10.4230/LIPICS.TYPES.2020.8 .
\bibitem[34]{source34} Jonas Kaiser, Steven Schäfer, and Kathrin Stark. Binder aware recursion over well-scoped de Bruijn syntax. In June Andronick and Amy P. Felty, editors, Proceedings of the 7th ACM SIGPLAN International Conference on Certified Programs and Proofs, CPP 2018, Los Angeles, CA, USA, January 8-9, 2018, pages 293-306. ACM, 2018. URL: https://doi.org/10.1145/3167098 .
\bibitem[35]{source35} Ambrus Kaposi. Formalisation of type checking into algebraic syntax. https://bitbucket.org/akaposi/tt-in-tt/src/master/Typecheck.agda , 2018.
\bibitem[36]{source36} Ambrus Kaposi. Quotient inductive-inductive types and higher friends. Talk given at the Homotopy Type Theory Electronic Seminar Talks (HoTTEST), October 2020. URL: https://akaposi.github.io/pres\_hottest.pdf .
\bibitem[37]{source37} Ambrus Kaposi, Simon Huber, and Christian Sattler. Gluing for type theory. In Herman Geuvers, editor, 4th International Conference on Formal Structures for Computation and Deduction, FSCD 2019, June 24-30, 2019, Dortmund, Germany, volume 131 of LIPIcs, pages 25:1-25:19. Schloss Dagstuhl - Leibniz-Zentrum für Informatik, 2019. URL: https://doi.org/10.4230/LIPICS.FSCD.2019.25 .
\bibitem[38]{source38} Ambrus Kaposi and András Kovács. A syntax for higher inductive-inductive types. In Hélène Kirchner, editor, 3rd International Conference on Formal Structures for Computation and Deduction, FSCD 2018, July 9-12, 2018, Oxford, UK, volume 108 of LIPIcs, pages 20:1-20:18. Schloss Dagstuhl - Leibniz-Zentrum für Informatik, 2018. URL: https://doi.org/10.4230/LIPICS.FSCD.2018.20 .
\bibitem[39]{source39} Ambrus Kaposi, András Kovács, and Thorsten Altenkirch. Constructing quotient inductive-inductive types. Proc. ACM Program. Lang., 3(POPL):2:1-2:24, 2019. URL: https://doi.org/10.1145/3290315 .
\bibitem[40]{source40} Ambrus Kaposi, András Kovács, and Ambroise Lafont. For finitary induction-induction, induction is enough. In Marc Bezem and Assia Mahboubi, editors, 25th International Conference on Types for Proofs and Programs, TYPES 2019, June 11-14, 2019, Oslo, Norway, volume 175 of LIPIcs, pages 6:1-6:30. Schloss Dagstuhl - Leibniz-Zentrum für Informatik, 2019. URL: https://doi.org/10.4230/LIPICS.TYPES.2019.6 .
\bibitem[41]{source41} Ambrus Kaposi and Szumi Xie. Type theory in type theory using single substitutions. In 30th International Conference on Types for Proofs and Programs. IT University of Copenhagen, 2024.
\bibitem[42]{source42} András Kovács and Ambrus Kaposi. Large and infinitary quotient inductive-inductive types. In Holger Hermanns, Lijun Zhang, Naoki Kobayashi, and Dale Miller, editors, LICS '20: 35th Annual ACM/IEEE Symposium on Logic in Computer Science, Saarbrücken, Germany, July 8-11, 2020, pages 648-661. ACM, 2020. URL: https://doi.org/10.1145/3373718.3394770 .
\bibitem[43]{source43} András Kovács. Type-Theoretic Signatures for Algebraic Theories and Inductive Types. PhD thesis, Eötvös Loránd University, Hungary, 2022. https://arxiv.org/abs/2302.08837.
\bibitem[44]{source44} Paul Blain Levy, John Power, and Hayo Thielecke. Modelling environments in call-by-value programming languages. Inf. Comput., 185(2):182-210, 2003. URL: https://doi.org/10.1016/S0890-5401(03)00088-9 .
\bibitem[45]{source45} Hugo Moeneclaey. Parametricity and semi-cubical types. In 36th Annual ACM/IEEE Symposium on Logic in Computer Science, LICS 2021, Rome, Italy, June 29 - July 2, 2021, pages 1-11. IEEE, 2021. URL: https://doi.org/10.1109/LICS52264.2021.9470728 .
\bibitem[46]{source46} Pierre-Marie Pédrot. Russian constructivism in a prefascist theory. In Holger Hermanns, Lijun Zhang, Naoki Kobayashi, and Dale Miller, editors, LICS '20: 35th Annual ACM/IEEE Symposium on Logic in Computer Science, Saarbrücken, Germany, July 8-11, 2020, pages 782-794. ACM, 2020. URL: https://doi.org/10.1145/3373718.3394740 .
\bibitem[47]{source47} Brigitte Pientka and Jana Dunfield. Beluga: A framework for programming and reasoning with deductive systems (system description). In Jürgen Giesl and Reiner Hähnle, editors, Automated Reasoning, pages 15-21, Berlin, Heidelberg, 2010. Springer Berlin Heidelberg.
\bibitem[48]{source48} Benjamin C. Pierce. Types and programming languages. MIT Press, 2002.
\bibitem[49]{source49} Benjamin C. Pierce, Arthur Azevedo de Amorim, Chris Casinghino, Marco Gaboardi, Michael Greenberg, Cătălin Hriţcu, Vilhelm Sjöberg, Andrew Tolmach, and Brent Yorgey. Programming Language Foundations, volume 2 of Software Foundations. Electronic textbook, 2024. Version 6.5, URL: http://softwarefoundations.cis.upenn.edu .
\bibitem[50]{source50} Jonathan Sterling. First Steps in Synthetic Tait Computability: The Objective Metatheory of Cubical Type Theory. PhD thesis, Carnegie Mellon University, USA, 2022. URL: https://doi.org/10.1184/R1/19632681.V1 .
\bibitem[51]{source51} Jonathan Sterling and Carlo Angiuli. Normalization for cubical type theory. In 36th Annual ACM/IEEE Symposium on Logic in Computer Science, LICS 2021, Rome, Italy, June 29 - July 2, 2021, pages 1-15. IEEE, 2021. URL: https://doi.org/10.1109/LICS52264.2021.9470719 .
\bibitem[52]{source52} Taichi Uemura. Abstract and Concrete Type Theories. PhD thesis, University of Amsterdam, 2021.
\bibitem[53]{source53} Philip Wadler, Wen Kokke, and Jeremy G. Siek. Programming Language Foundations in Agda. Online, August 2022. URL: https://plfa.inf.ed.ac.uk/22.08/ .
\end{thebibliography}
\end{document}
